\cotefile{136}
\scanpage{95}\makeatletter\tagsleft@true\makeatother\mspage{- 169 -}\hfill IX

finies ; donc $X\longmapsto\hat X$ transforme groupes en groupes. Notons que $\hat X$ ne dépend en fait que de la restriction de $X$ à un voisinage ouvert arbitraire $U$ de $S$.

\uline{Lemme} 13.2.1. \uline{Les données sont celles de} 13.1, \uline{et on suppose que les conditions suivantes sont vérifiées} : a) $A=A^{\unsure{o|\circ}}$, b) $R$ \uline{est constant au voisinage de} $s$, \uline{et} c) $R_s$ \uline{engendre l'espace vectoriel complexe} $E_s$. \uline{Soit} $A'$ \uline{un deuxième groupe analytique lisse sur} $S$, \uline{et considérons l'application}
\[
 u\longmapsto\widehat{\unsure{\mathbb{u}|u}}:\operatorname{Hom}_{S\text{-}\mathrm{gr}}(A,A')\longrightarrow\operatorname{Hom}_{\hat S\text{-}\mathrm{gr}}(\widehat{\unsure{\mathbb A|A}},\hat A')\quad,
\]
\uline{où les notations sont celles de} 13.2, \uline{et l'application analogue obtenue en se restreignant au-dessus d'un voisinage ouvert} $U$ \uline{variable de} $s$, \uline{et passant à la limite inductive} :
\[
 u\longmapsto\hat u:\operatorname{Hom}_{(S,s)\text{-}\mathrm{gr}}(A,A')\longrightarrow\operatorname{Hom}_{\hat S\text{-}\mathrm{gr}}(\hat A,\hat A')\quad,
 \tag*{(13.2.1.1)}
\]
(\uline{où} $(S,s)$ \uline{désigne le} germe \uline{d'espace analytique de} $S$ \uline{en} $s$). \uline{Alors l'application précédente est} bijective.

Considérons en effet la suite exacte
\[
 0\longrightarrow R'\longrightarrow E'\longrightarrow {A'}^{o}\longrightarrow\unsure{0|C}
\]
analogue à (13.1.3) définie par $A'$. Compte tenu du fait que les homomorphismes $A\longrightarrow A'$ se factorisent nécessairement par $ {A'}^{o}$, et de même après les changements de base $U\longrightarrow S$ et $S_n\longrightarrow S$, on peut supposer $A'={A'}^{o}$. Compte tenu du dictionnaire 13.1.4, il faut prouver que la donnée d'un
% archive sheet number: 94
% stamp: numbered label 487; 487 also in pencil above the lower strip.
